% (c) Nikita Lisitsa, lisyarus@gmail.com, 2025

\documentclass[10pt]{beamer}

\usepackage[T2A]{fontenc}
\usepackage[russian]{babel}
\usepackage{minted}

\usepackage[most]{tcolorbox}
\tcbuselibrary{minted}

\usepackage{graphicx}
\graphicspath{ {./images/} }

\usetheme{metropolis}

\usepackage{adjustbox}

\usepackage{color}
\usepackage{soul}

\usepackage{hyperref}

\definecolor{codebg}{RGB}{29,35,49}

\setminted{fontsize=\footnotesize}

\setlength\partopsep{-\topsep}

\definecolor{blue}{rgb}{0,0,1}
\definecolor{red}{rgb}{1,0,0}

\makeatletter
\newcommand{\slideimage}[1]{
  \begin{figure}
    \begin{adjustbox}{width=\textwidth, totalheight=\textheight-2\baselineskip-2\baselineskip,keepaspectratio}
      \includegraphics{#1}
    \end{adjustbox}
  \end{figure}
}
\makeatother

\title{Компьютерная графика}
\subtitle{Практика 3: Буферы, атрибуты вершин, кривые Безье}
\date{2026}

\setbeamertemplate{footline}[frame number]

\begin{document}

\frame{\titlepage}

\begin{frame}[fragile]
\frametitle{Кривые Безье}
\begin{itemize}
\item Используются для генерации плавных кривых
\pause
\item Строются по набору точек \begin{math}p_0, p_1, \dots, p_n\end{math}
\pause
\item Кривая Безье с параметром \begin{math}t \in [0, 1]\end{math} определяется как \textit{аффинная комбинация}

\begin{center}
\begin{math}
b(t) = \sum\limits_{k=0}^n b_{k,n}(t) \cdot p_k
\end{math}
\end{center}

\pause
\item Коэффициенты -- полиномы Бернштейна:

\begin{center}
\begin{math}
b_{k,n}(t) = \binom{n}{k}t^k(1-t)^{n-k}
\end{math}
\end{center}

\pause
\item Кривая первого порядка (\begin{math}n=1\end{math}): отрезок \begin{math}p_0 p_1\end{math}

\begin{center}
\begin{math}
b(t) = (1-t)\cdot p_0 + t \cdot p_1
\end{math}
\end{center}

\pause
\item Кривая второго порядка (\begin{math}n=2\end{math}):

\begin{center}
\begin{math}
b(t) = (1-t)^2 \cdot p_0 + 2t(1-t) \cdot p_1 + t^2 \cdot p_2
\end{math}
\end{center}

\end{itemize}
\end{frame}

\begin{frame}[fragile]
\frametitle{Задание 1}
Загружаем данные в буфер на GPU
\usemintedstyle{solarized-light}
\begin{itemize}
\item Создайте буфер: \mintinline{cpp}|device.create_buffer|, \mintinline{cpp}$usage = CopyDst | Vertex$, размер -- 3 вершины (надо посчитать в байтах!)
\item Загрузите в него данные \mintinline{cpp}|queue.write_buffer|
\end{itemize}
\end{frame}

\begin{frame}[fragile]
\frametitle{Задание 2}
Настраиваем атрибуты
\usemintedstyle{solarized-light}
\begin{itemize}
\item Настройте атрибуты вершин при создании пайплайна (offset, stride, и формат надо понять по структуре \mintinline{cpp}|vertex|)
\item Перед \mintinline{cpp}|render_pass.draw| выставьте наш буфер как вершинный \mintinline{cpp}|render_pass.set_vertex_buffer| -- должен нарисоваться треугольник
\end{itemize}
\end{frame}

\begin{frame}
\frametitle{Задание 2}
\slideimage{2.png}
\end{frame}

\begin{frame}[fragile]
\frametitle{Задание 3}
Переходим к оконным координатам
\usemintedstyle{solarized-light}
\begin{itemize}
\item Заполните массив \mintinline{cpp}|viewMatrix| матрицей, которае совершает проеобразование

\begin{center}
\begin{math}
X: [0, width] \rightarrow [-1, 1]
\end{math}

\begin{math}
Y: [height, 0] \rightarrow [-1, 1]
\end{math}
\end{center}

\item Измените координаты треугольника, чтобы он был заметен на экране
\end{itemize}
\end{frame}

\begin{frame}
\frametitle{Задание 3}
\slideimage{3.png}
\end{frame}

\begin{frame}[fragile]
\frametitle{Задание 4}
Динамически добавляем/удаляем точки мышкой
\usemintedstyle{solarized-light}
\begin{itemize}
\item Сделайте массив вершин изначально пустым
\item При нажатии левой кнопки мыши (\mintinline{cpp}|SDL_BUTTON_LEFT|) добавьте новую вершину в массив с координатами мыши (цвета выбирайте как угодно)
\item При нажатии правой кнопки мыши (\mintinline{cpp}|SDL_BUTTON_RIGHT|) удалите последнюю вершину из массива, если он не пустой
\item Если массив с вершинами изменился, обновите данные в буфере (возможно, его придётся пересоздать с большим размером)
\item Меняем тип примитива у пайплайна на \mintinline{cpp}|LineStrip|
\item Меняем количество вершин в \mintinline{rust}|render_pass.draw()| на размер массива вершин
\item Данные должны обновляться \textbf{\alert{только}} когда линия изменилась!
\end{itemize}
\end{frame}

\begin{frame}
\frametitle{Задание 4}
\slideimage{4.png}
\end{frame}

\begin{frame}[fragile]
\frametitle{Задание 5}
Генерируем и рисуем кривые Безье
\usemintedstyle{solarized-light}
\begin{itemize}
\item Создаёте ещё один буфер и массив вершин для точек кривой Безье
\item Заведите параметр, управляющий уровнем детализации кривой: \mintinline{cpp}|int quality = 4;|
\item При добавлении/удалении точек пересчитываем \textit{заново} всю кривую Безье
\item Если в исходной ломаной \mintinline{cpp}|N| \textit{отрезков}, то в кривой Безье должно быть \mintinline{cpp}|N * quality| \textit{отрезков}
\item Параметр \mintinline{cpp}|t| должен равномерно меняться от 0 до 1 вдоль всей кривой
\item Цвет -- любой, но отличающийся от цвета исходной ломаной (чтобы её было лучше видно)
\item Данные в буферах должны обновляться \textbf{\alert{только}} при их изменении
\item Рисуем кривую, используя тот же пайплайн (но другой буфер)
\end{itemize}
\end{frame}

\begin{frame}
\frametitle{Задание 5}
\slideimage{5.png}
\end{frame}

\begin{frame}[fragile]
\frametitle{Задание 6}
Динамически меняем уровень детализации
\usemintedstyle{solarized-light}
\begin{itemize}
\item При нажатии на стрелку влево (\mintinline{cpp}|SDL_LEFT|) уровень детализации \mintinline{cpp}|quality| должен уменьшаться (но не может быть меньше 1)
\item При нажатии на стрелку вправо (\mintinline{cpp}|SDL_RIGHT|) уровень детализации \mintinline{cpp}|quality| должен увеличиваться
\item Данные в буферах должны обновляться \textbf{\alert{только}} при их изменении!
\end{itemize}
\end{frame}

\begin{frame}
\frametitle{Задание 6}
\slideimage{6.png}
\end{frame}

\begin{frame}[fragile]
\frametitle{Задание 7*}
\usemintedstyle{solarized-light}
Добавляем ползающий пунктир к кривой Безье
\begin{itemize}
\item Пунктира можно добиться, не рисуя половину пикселей кривой в зависимости от расстояния до начала
\item Например, пиксели кривой на расстоянии \mintinline{cpp}|[0, 20]| от начала рисуются, на расстоянии \mintinline{cpp}|[20, 40]| от начала не рисуются, \mintinline{cpp}|[40, 60]| рисуются, и т.д.
\item Для этого нужно добавить в вершины кривой расстояние вдоль кривой от её начала
\item Расстояние до первой вершины = 0
\item Расстояние до  вершины \mintinline{cpp}|i| = расстояние до  вершины \mintinline{cpp}|(i-1)| + расстояние от вершины \mintinline{cpp}|(i-1)| до вершины \mintinline{cpp}|i|
\item Для расстояния между вершинами пригодится функция \mintinline{cpp}|math::length|
\item Расстояние нужно интерполировать между вершинами и передать во фрагментный шейдер
\end{itemize}
\end{frame}

\begin{frame}[fragile]
\frametitle{Задание 7*}
\usemintedstyle{solarized-light}
\begin{itemize}
\item Чтобы понять, нужно ли рисовать пиксель, пригодится функция WGSL \mintinline{wgsl}|modf|
\item Чтобы отменить рисование пикселя, нужно во фрагментном шейдере вызвать команду \mintinline{wgsl}|discard;|
\item Само расстояние нужно добавить как поле вершины и как атрибут в пайплайне и в вершинном шейдере
\item Чтобы пунктир двигался, можно прибавить к расстоянию что-то, зависящее от времени (придётся передать время в шейдер как часть immediates)
\item Чтобы на исходную ломаную не применялся пунктир, придётся передать ещё какой-нибудь флаг в immediates
\end{itemize}
\end{frame}

\begin{frame}
\frametitle{Задание 7*}
\slideimage{7.png}
\end{frame}

\end{document}
